\chapter{Magnetic Flux Tubes}
\label{ch:fluxtubes}
\acresetall


\begin{summary}
This chapter serves as an introduction to the work presented in chapters 
\ref{ch:greensfunc} through \ref{ch:periodic}.  Neutron stars have a crust made 
of dense nuclear material which is generally expected to be in a type-II superconducting state. 
We therefore expect the intense magnetic fields in the crust of a neutron star 
to form a lattice of individual flux tubes. This provides the primary motivation for the 
following chapters which study magnetic flux tubes in \ac{QED}. In this chapter we will 
review some of the basic properties of the dense nuclear material of neutron stars that 
may result in the formation of a flux tube lattice.
\end{summary}



\section{Introduction}
In this part, we will explore the \ac{QED} effective action 
in cylindrically symmetric, extended tubes of magnetic flux.
These configurations may be called 
flux tubes, strings, or vortices, depending on the context.  
Flux tubes are of interest in astrophysics 
because they describe magnetic structures near stars and planets, 
cosmic strings~\cite{vilenkin2000cosmic}, and vortices in the superconducting core of neutron stars
\cite{2006pfsb.book..135S, schmitt2010dense}.
Outside of astrophysics, magnetic vortex systems are at the forefront of research in condensed 
matter physics for the role they play in superconducting systems and in \ac{QCD} research 
for their relation to center vortices, a gluonic configuration analogous to magnetic vortices 
which is believed to be important to quark confinement~\cite{tHooft19781, 2003PrPNP..51....1G}.  
In this part of the dissertation, we will primarily discuss the 
roles played by magnetic flux tubes in neutron star cores. 

Our motivation for discussing flux tubes comes from the fact that 
superconductivity is predicted in the nuclear matter of neutron stars
and that some superconducting materials produce a lattice of flux tubes 
when placed in an external magnetic field. Superconductivity is a macroscopic 
quantum state of a fluid of fermions that, most notably,
allows for the resistanceless conduction of charge. In 1933, 
Meissner and Ochsenfeld observed that magnetic fields are repelled 
from superconducting materials~\cite{meisner33}. 
In 1935, F. and H. London described the Meissner effect in terms of 
a minimization of the free energy of the superconducting current~\cite{london35}.
Then, in 1957, by studying the superconducting 
electromagnetic equations of motion in cylindrical coordinates, 
Abrikosov predicted the possible existence of line defects in superconductors 
which can carry quantized magnetic flux through the superconducting material
\cite{abrikosov57}. 

A more complete microscopic description of superconducting materials is given by 
BCS (Bardeen, Cooper, and Schrieffer) theory~\cite{PhysRev.108.1175}. 
However, detailed discussions 
are outside the scope of this chapter. Instead, I will simply outline the main 
features of superconductivity that motivate the study of flux tubes in neutron 
stars. Interested readers may pursue more thorough reviews of superconductivity 
and superfluidity in neutron stars~\cite{2006pfsb.book..135S, schmitt2010dense}.

\section{Superconductivity}


To see how superconductivity is possible, 
consider fermions at zero temperature with chemical potential, $\mu$. The
free energy of this system is 
\be
	\Omega = E - \mu N.
\ee
For non-interacting fermions, we could add fermions at the Fermi surface 
without changing the free energy, since the first and second terms would 
change by the same amount. If, instead, the fermions are interacting, 
the binding energy between the fermions means that the free energy can 
be reduced by adding fermions. So, if there is an attractive interaction 
between the fermions, then there is a new ground state of the fermion fluid
where pairs of fermions are created at the Fermi surface. These pairs of 
fermions are called Cooper pairs and (since a pair of fermions can be viewed 
as a boson) they form a Bose condensate.

The particles in a Cooper pair do not form a bound state, but they are interacting 
through attractive forces, so there is an energy $\Delta$ associated with 
separating them. The consequence of this is that there is an energy gap in 
the dispersion relation for the Cooper pair:

\be
	\epsilon_k = \sqrt{(\sqrt{k^2+m^2} - \mu)^2 + \Delta^2}.
\ee
This energy gap means that a finite amount of energy is required to 
excite a single electron state, even near the Fermi surface. 
An amazing consequence of this is that particles flowing through 
the Bose condensate cannot scatter inelastically 
from the phonons because 
fermions cannot be excited at low energies. For neutral fermions, 
such as neutrons, this can give rise to superfluidity, characterized by 
frictionless flows. If the fermions forming 
the cooper pairs are charge carriers such as protons or electrons, 
this effect gives rise to superconductivity, characterized by the 
resistanceless conduction of charge. 

\section{Energy Interpretation of the \texorpdfstring{\acs{QED}}{QED} Effective Action}

In section \ref{sec:qedea}, I argued that $E[J^\mu, \bar{\eta},\eta]$ was the 
vacuum energy associated with the external source. Thus, equation (\ref{eqn:LegendreTrans})
implies that the effective action is closely related to the vacuum energy of a 
classical field configuration. Importantly, the effective action 
can be thought of as the additive inverse of the energy. 
Since this idea is important for understanding 
the free energy associated with a magnetic flux tube, I will make the relationships 
between the Hamiltonian expectation value, $E[J^\mu, \bar{\eta},\eta]$, 
and $\Gamma[A^0_\mu, \bar{\psi}^0, \psi^0]$ clear and explicit
\footnote{This section follows section 16.3 of Weinberg~\cite{weinberg1966quantum}.}.

To study the vacuum energy of a static electromagnetic field, we would like to 
minimize the expectation value of the Hamiltonian, 
\be
	\mean{H}_\Omega = \bra{\Omega}H\ket{\Omega},
\ee
under the constraint that the vacuum expectation value of the quantum gauge 
field is the classical field,
\be
	\bra{\Omega}A_\mu(\vec{x})\ket{\Omega} = A_\mu^0(\vec{x}).
\ee
We assume that the classical fields, $\bar{\psi}^0$ and $\psi^0$, vanish.
Using the method of Lagrange multipliers, the function which we would like to 
minimize is:
\be
	\mean{H}_\Omega - \alpha\braket{\Omega}{\Omega}
	-\int d^3x \beta^\mu(\vec{x})\bra{\Omega}A_\mu(\vec{x})\ket{\Omega}.
\ee
So, we find that we must satisfy
\be
	\label{eqn:Hmin}
	H\ket{\Omega} = \alpha\ket{\Omega} + \int d^3x 
	\beta^\mu(\vec{x}) A_\mu(\vec{x}) \ket{\Omega}.
\ee
This can be done if the Lagrange multipliers 
$\alpha$ and $\beta^\mu(\vec{x})$ are functionals of $A_\mu^0(\vec{x})$.

Let $\ket{\Psi}_{J^\mu}$ be normalized eigenvectors of the Hamiltonian in the presence of 
external source $J^\mu$. The energy eigenvalue equation is 

\be
	\left[H-\int d^3x\left(J^\mu A_\mu + \bar{\eta}\psi + \bar{\psi}\eta\right)\right]
	\ket{\Psi}_{J^\mu} = \frac{E[J^\mu,\bar{\eta},\eta]}{\mathcal{T}}\ket{\Psi}_{J^\mu},
\ee
where $\mathcal{T}$ is the time extent of the functional integration.
If we turn on the external source $J^{0,\mu}$ adiabatically to put the vacuum into the 
energy eigenstate $\ket{\Psi}_{J^{0,\mu}}$, then equation (\ref{eqn:Hmin}) is satisfied by 
taking 
\be
	\ket{\Omega} = \ket{\Psi}_{J^{0,\mu}},
\ee
\be
	\alpha = E[J^{0, \mu}, 0, 0],
\ee
and
\be
	\beta_\mu(\vec{x}) = J_\mu^0(\vec{x}).
\ee
Under these substitutions, equation (\ref{eqn:Hmin}) provides an expression for the 
vacuum expectation value for the Hamiltonian in the external field in terms of the 
external fields and currents.
\be
	\mean{H}_{A_\mu^0} = \frac{E[J_\mu^0,0,0]}{\mathcal{T}} + \frac{\int d^4x J_\mu^0A^{0,\mu}}{\mathcal{T}}
\ee
The above equation is related simply to the definition of the effective action 
as a Legendre transformation of the energy
\be
	\tag{\ref{eqn:LegendreTrans}}
	\Gamma[A_\mu^0, \bar{\psi}^0, \psi^0] = -E[J^\mu, \bar{\eta},\eta]
	-\int d^4 y (J^\mu(y)A_\mu^0(y) + \bar{\eta}(y)\psi^0(y) +\bar{\psi}^0(y)\eta(y)).
\ee
So, we arrive at the relationship between the expectation value of the 
Hamiltonian in the external field and the effective action.
\be
	\mean{H}_{A_\mu^0} = -\frac{1}{\mathcal{T}} \Gamma[A_\mu^0]
\ee
Because the energy is extensive, the most interesting quantity for magnetic 
flux tube configurations is the total energy per unit length.
\be
	\frac{\mean{H}_{A_\mu^0}}{L_z} = -\frac{\Gamma[A_\mu^0]}{L_z \mathcal{T}}
\ee
Thus, we have explicitly confirmed the expectation that the negative effective 
action is related to the total vacuum energy of a specific field configuration.

\section{Flux Tube Free Energy}

Many of the basic properties of flux tubes in superconductors can 
be understood by examining the free energy in Ginzberg-Landau theory.
The Ginzberg-Landau free energy is given by


\ba
	\label{eqn:GLFE}
	E &=& \int d^2 x \biggr\{\frac{1}{2m}\biggr|\left(\vec{\nabla }
		- iq\vec{A}(\vec{x})\right)\psi(\vec{x})\biggr|^2 \nonumber \\
	 & & - \mu | \psi(\vec{x}) |^2+\frac{a}{2}| \psi(\vec{x}) |^4 
	 + \frac{1}{8\pi}(\vec{\nabla} \times \vec{A} (\vec{x}))^2\biggr\},	
\ea
where $\psi(\vec{x})$ is a complex order field such that the 
density of superconducting fermions is proportional to $|\psi|^2$, $\mu$ 
is the chemical potential and $a$ is related to the scattering length. 
The first term represents the dynamical energy of fluctuations of the 
order field. The next two terms represent the potential energy 
of the order field. The final term is the familiar energy of the 
classical magnetic field.

Incorporating the 1-loop effects from \ac{QED}, the classical term 
must be replaced with the energy of the field, 
$E_{\rm 1-loop} = -\frac{\Gamma[A_\mu^0]}{L_z\mathcal{T}}$.
\ba
	E &=& \int d^2 x \biggr\{\frac{1}{2m}\left|\left(\vec{\nabla }
		- iq\vec{A}(\vec{x})\right)\psi(\vec{x})\right|^2 \nonumber \\
	 & & - \mu|\psi(\vec{x})| ^2+\frac{a}{2}|\psi(\vec{x})|^4 
	 \biggr\} -\frac{\Gamma[A_\mu^0]}{\mathcal{T}}.	
\ea
This addition can be viewed as a correction to the classical 
magnetic field energy in the Ginzburg-Landau free energy.
Many important behaviours of superconductors are determined 
by minimizing the free energy. The technique discussed in section 
\ref{sec:stationaryEA} may lead to a way to perform
this minimization while including 1-loop \ac{QED} effects.

The Ginzburg-Landau free energy plays an important role in 
determining the difference between type-I and type-II superconductors.
For example, the free energy of a pair of flux tubes will contain 
terms characterizing the energy of each flux tube, but also 
cross terms arising from the interaction between these flux tubes.
The sign of this interaction energy determines if the flux tubes 
will repel each other and form a lattice, or attract each other and collapse.
In the classical case, when the fields are separated spatially, there 
is no interaction energy arising from the last term in equation (\ref{eqn:GLFE}).
However, when the 1-loop \ac{QED} effects are taken into account, the 
non-local effects arising due to a nearby flux tube make a contribution 
to the free energy~\cite{Langfeld:2002vy}. 

\subsection{Meissner Effect}

The free energy in a superconductor is minimized if the magnetic 
field obeys its equation of motion, the London equation,

\be
	\vec{\nabla} \times \vec{\nabla} \times \vec{B}(\vec{x})
	= - \lambda_L^2\vec{B}(\vec{x}),
\ee
where $\lambda_L$ is called the London penetration depth,
which is defined as
\be
	\label{eqn:LondonPD}
	\lambda_L = \sqrt{\frac{2m_f c^2}{4\pi (2 q_f)^2 n_0}}
\ee
where $m_f$ and $q_f$ are the fermion mass and charge, 
respectively, and $n_0$ is the fermion density.

The solution to this equation implies that magnetic 
fields are repelled exponentially from the surface
of the superconductor.

\be
	B(x) = e^{-\frac{x}{\lambda_L}}
\ee


Superconducting regions of a material must repel magnetic 
fields according to this Meissner effect. However, there 
can be a large free energy cost to expelling field lines 
around a superconducting region. For some materials, it 
is therefore energetically favourable to allow magnetic field 
lines to penetrate the material through a non-superconducting 
core. When this occurs, tubes of magnetic flux are 
permitted in the material.

\subsection{Interaction Energy Between Vortices}


The interaction energy between two vortices, and therefore 
the force between them, can be found by examining the 
free energy of a two vortex system and subtracting off 
the energy expected from each individual flux tube~\cite{PhysRevB.3.3821}.
The interaction energy is then given by the remaining cross terms

\be
	E_{\rm int} = \int d^2 x \left(J_1^\mu A_\mu^2 - \rho_1 \psi_2\right)
\ee
where $\psi_i$ and $A_\mu^i$ are the fields induced by sources 
$\rho_i$ and $J_i^\mu$~\cite{Speight_staticintervortex}.

In a superconductor, these two terms tend to oppose one 
another. Two neighbouring vortices 
have charged fermion currents in opposite directions 
and are repelled from each other. However, there 
is also an attractive force due to a free energy associated 
with defects in the superconductor which can be reduced 
if two vortices are combined~\cite{PhysRevB.3.3821}. 
So, the vortices can be 
either attractive or repulsive depending on the relative 
sizes of these terms.


This leads to two different behaviours for superconductors 
in the presence of magnetic fields. The behaviour is 
naively determined by the Ginzburg-Landau parameter,

\be
	\label{eqn:GLparam}
	\kappa = \frac{\lambda_L}{\xi}.
\ee

The coherence length, $\xi$, is the length scale associated with density variations of superconducting fermions. This
length scale is given in terms of $a$ in the Ginzburg-Landau theory (see equation (\ref{eqn:GLFE})), 
or in terms of the energy gap, $\Delta$, and the velocity of Cooper pairs, $v_f$ in BCS theory. 
\ba
	\xi &=& \sqrt \frac{\hbar^2}{4 m_f |a|} \\
	 &=& \frac{2 \hbar v_f}{\pi \Delta}
\ea
When $\kappa > \frac{1}{\sqrt{2}}$, the vortices will 
repel, and when $\kappa < \frac{1}{\sqrt{2}}$, the vortices 
will attract. Of course, this analysis has assumed that 
the free energy has been completely accounted for by 
the standard analysis. If there are additional 
considerations in the interaction energy, they will 
influence the properties of the superconducting material.
Previously, the effects of an asymmetry in the scattering 
length between neutron and proton cooper 
pairs~\cite{2004PhRvL..92o1102B},  and the 
effects of currents transported along vortices~\cite{2007PhRvC..76a5801C} have been 
suggested as mechanisms for making a type-II superconductor 
behave like a type-I superconductor. Friction between superconducting and 
normal fluid domains in a type-I proton superfluid 
causes a dissipation which is consistent with precession~\cite{2005PhRvD..71h3003S}.
Thus, a major motivation 
for studying the \ac{QED} contribution to the free energy 
of flux tubes, is to learn how this contribution may affect their 
interaction energies and the rotational dynamics of neutron stars.

If $\kappa$ is large, meaning the vortices repel, 
the superconductor can be in a so-called mixed state where the 
vortices will form a triangular Abrikosov lattice
with each vortex carrying a single quanta of flux~\cite{abrikosov57}. 
The superconductors where this happens allow partial 
penetration of the magnetic field and are called 
type-II superconductors. However, this mixed state only occurs for 
magnetic fields strong enough to form vortices, but not so strong that they 
destroy superconductivity. If $\kappa$ is small, 
the vortices will attract each other and annihilate. 
When this happens, the superconductor does not allow 
magnetic flux to penetrate the material and is called 
a type-I superconductor.

\begin{figure}[h]
	\centering
		\includegraphics[width=8.5cm]{images/hexlatticegs.eps}
                \caption[Hexagonal lattice]
				{In type-II superconductors, flux tubes 
				form an Abrikosov lattice. The geometry is 
				shown in the diagram above where each 
				hexagonal space can be imagined as containing 
				a flux tube at its centre}
	\label{fig:hexlatticegs}
\end{figure}

In this section, I have argued that without the electromagnetic 
forces, the vortices in a superconductor will experience an 
attractive force. Naively, this suggests that superfluids 
should be type-I since they lack a coupling to the gauge fields.
This contrasts with the experimental observation that 
superfluids are type-II. The subtlety here is that the 
nature of the self-interaction of the order field is different 
in the absence of gauge fields, and is always repulsive for 
superfluids. We can understand a superconductor as the $q_e \to 0$
limit of a superfluid. In this limit, the London penetration depth, equation 
(\ref{eqn:LondonPD}), becomes infinite and so does the Ginzburg-Landau
parameter, equation (\ref{eqn:GLparam}).
\be
	\lim_{q_e\to0}\kappa = \lim_{q_e\to 0}\frac{\lambda_L}{\xi} 
	\sim \lim_{q_e\to0}\frac{1}{q_e} = \infty
\ee
So, we predict type-II behaviour from superfluids, consistent with 
experimental results. The self-interaction of the order field 
is repulsive in the absence of a gauge field.

\section{Nuclear Superconductivity in Neutron Stars}

In the dense nuclear matter of a neutron star, it may be possible 
to have neutron superfluidity, proton superconductivity, and 
even quark colour superconductivity~\cite{2006pfsb.book..135S, schmitt2010dense}. 
The first prediction of 
neutron-star superfluidity dates back to Migdal in 1959~\cite{1959NucPh..13..655M}.
The arguments which make superfluidity seem likely are based on 
the temperatures of neutron stars. A short time after their 
creation, neutron stars are very cold. The temperature in the 
interior may be a few hundred keV. Studies of nuclear matter 
show that the transition temperature is 
$T_c \gtrsim 500$ keV~\cite{1989tns..conf..457S}.
So, it is expected that the nuclear matter in a neutron star 
forms condensates of cooper pairs.

Because of this, some fraction of neutrons in the inner crust 
of a neutron star are expected to be superfluid. These neutrons make up 
about a percent of the moment of inertia of the star and are 
weakly coupled to the nuclear crystal lattice which makes up the 
remainder of the inner crust. If the neutron vortices in the 
superfluid component move at nearly the same speed as the nuclear 
lattice, the vortices can become pinned to the lattice so that the 
superfluid shares an angular velocity with the crust. 
This pinning between the fluid and solid crust has observable impacts 
on the rotational dynamics of the neutron star, for the same reasons that 
hard-boiling an egg (pinning the yolk to the shell) produces an 
observable difference in the way it spins on its side.

This picture of a neutron superfluid co-rotating with a solid crust 
has been used to interpret several types of pulsar timing anomalies.
Pulsars are nearly perfect clocks, although they gradually spin-down 
as they radiate energy. Occasionally, though, pulsars demonstrate 
deviations from their expected regularity. A glitch is an abrupt 
increase in the rotation and spin-down rate of a pulsar, followed 
by a slow relaxation to pre-glitch values over weeks or years. This 
behaviour is consistent with the neutron superfluid suddenly 
becoming unpinned from the crust and then dynamically relaxing 
due to its weak coupling to the crust until it is pinned once 
again~\cite{1975Natur.256...25A}.

The neutron star in Cassiopeia A has been observed to be rapidly cooling
\cite{2010ApJ...719L.167H}. The surface temperature decreases by about 
4\% every 10 years. This observation is also strong evidence 
of superfluidity and superconductivity in neutron stars
\cite{2011PhRvL.106h1101P, 2011MNRAS.412L.108S}. The observed cooling 
is too fast to be explained by the observed x-ray emissions 
and standard neutrino cooling. However, 
the cooling is readily explained by the emission of neutrinos during 
the formation of neutron Cooper pairs. Based on such a model, 
the superfluid transition temperature of neutron star matter is
$\sim 10^{9}$ K or $\sim 90$ keV. 

Further hints regarding superfluidity in neutron stars 
come from long-term periodic variability
in pulsar timing data. For example, variabilities in PSR B1828-11
were initially interpreted as free 
precession (or wobble) of the star~\cite{2000Natur.406..484S}. 
If neutron stars can
precess, observations could strongly constrain the ratio of the moments
of inertia of the crust and the superfluid neutrons. Moreover, the existence
of flux tubes (\ie type-II superconductivity)
in the crust is generally incompatible with the slow, large 
amplitude precession suggested by PSR B1828-11~\cite{2000Natur.406..484S}. The neutron vortices 
would have to pass through the flux tubes, which should cause a huge dissipation 
of energy and a dampening of the precession which is not observed~\cite{PhysRevLett.91.101101}. 
However,
recent arguments suggest that the timing variability data is not well explained 
by free precession and that it more likely suggests that the star is switching between 
two magnetospheric states~\cite{2010Sci...329..408L}. 
Nevertheless, other authors suggest not being premature in throwing 
out the precession hypothesis without further observations
\cite{2012MNRAS.420.2325J}.

\section{Magnetic Flux Tubes in Neutron Stars}


If a magnetic field is able to penetrate the proton superfluid on a microscopic level,
it must do so by forming a triangular Abrikosov 
lattice with a single quanta of flux in each flux tube. So, the 
density of flux tubes is given simply by the average field strength.
If the distance between flux tubes is $l_f$, the flux in a circular 
region within $l_f/2$ of a flux tube is given by 

\be
	F = \frac{2\pi \mathcal{F}}{e} = 2\pi \int_0^{l_f/2} B \rho d\rho
\ee
where we have introduced a dimensionless measure of flux 
$\mathcal{F} = \frac{e}{2\pi}F$.
So, the distance between flux tubes is
\be
	l_f = \sqrt{\frac{8 \mathcal{F}}{eB}}.
\ee
If the magnetic field is the critical field strength, $B_k$, then 
the flux tubes are separated by a few Compton electron wavelengths.
This is particularly interesting since this is the distance scale 
associated with non-locality in \ac{QED}.

The size of a flux tube profile in laboratory superconductors is 
on the nanometer or micron scale~\cite{poole2007superconductivity}. Because the flux 
is fixed, the size of the tube profile determines the strength of the magnetic field 
within the tube. For laboratory superconductors
the field strengths are small compared to the quantum critical field, and the field is slowly 
varying on the scale of the Compton wavelength.  In this case, the quantum corrections to the 
free energy are known to be much smaller than the classical contribution 
(see section \ref{sec:EAisoflux}).
The size scale for the flux tubes in a superconductor is determined by the London penetration depth. 
In a neutron star, this quantity is estimated to be a small fraction of a 
Compton wavelength, much smaller than in laboratory superconductors
\cite{lrr-2008-10, PhysRevLett.91.101101}. In this case, 
the magnetic field strength at the centre of the 
tube exceeds the quantum critical field strength and the field varies rapidly, 
rendering the derivative expansion description of the effective action unreliable.

The Ginzburg-Landau parameter, equation (\ref{eqn:GLparam}), is the 
ratio of the proton coherence length, $\xi_p \sim 30$ fm, and the 
London penetration depth of a proton super conductor, $\lambda_p \sim 80$~fm
\cite{PhysRevLett.91.101101}. 

\be
	\kappa = \frac{\lambda_p}{\xi_p} \sim 2
\ee
We therefore expect that the proton cooper pairs most likely 
form a type-II superconductor~\cite{Sedrakian:2006xm}. 
However, it is possible that 
physics beyond what is taken into account in the standard 
picture affects the free energy of a magnetic flux tube. In 
that case, the interaction between two flux tubes may indeed 
be attractive in which case the neutron star would be a type-I 
superconductor.

\section{\texorpdfstring{\acs{QED}}{QED} Effective Actions of Flux Tubes}
\label{sec:EAisoflux}

Vortices of magnetic flux have very important impacts on the quantum 
mechanics of electrons. In particular, the phase of the electron's wavefunction 
is not unique in such a magnetic field. This is demonstrated by the Aharonov-Bohm
effect~\cite{0370-1301-62-1-303,PhysRev.115.485}. The first calculations of the 
fermion effective energies of these configurations were for infinitely thin Aharonov-Bohm
flux strings~\cite{Gornicki1990271, 1998MPLA...13..379S}. Calculations for thin strings 
were also performed for cosmic string configurations~\cite{0264-9381-12-5-013}. For these 
infinitely-thin string magnetic fields, the energy density is singular for small radii. So, 
it is not possible to define a total energy per unit length. Another approach was to 
compute the effective action for a finite-radius flux tube where the magnetic flux was 
confined entirely to the radius of the tube~\cite{PhysRevD.51.810}. This approach 
results in infinite classical energy densities.

Physical flux tube configurations would have a finite radius.
The earliest paper to deal with finite radius flux tubes in QED considered the effective 
action of a step-function profiled flux tube using the Jost function of the related 
scattering problem~\cite{1999PhRvD..60j5019B}. One of the conclusions from this research 
was that the quantum correction to the classical energy was relatively 
small for any value of the flux tube size, for the entire range of applicability of \ac{QED}. 
The techniques from this study were 
soon generalized to other field profiles including, a delta-function cylindrical shell magnetic 
field~\cite{PhysRevD.62.085024}, and more realistic flux tube configurations 
such as the Gaussian~\cite{2001PhRvD..64j5011P} and the Nielsen-Olesen 
vortex~\cite{2003PhRvD..68f5026B}. 
Flux tube vacuum energies were also analyzed extensively using a spectral method
\cite{2005NuPhB.707..233G, 2006JPhA...39.6799W, Weigel:2010pf}.

The effective actions of flux tubes have been previously analyzed using worldline 
numerics~\cite{Langfeld:2002vy}.  This research investigated isolated 
flux tubes, but also made use of the loop cloud method's applicability to situations 
of low symmetry to investigate pairs of interacting 
vortices. One conclusion from that investigation was that the fermionic effects resulted 
in an attractive force between vortices with parallel orientations, and a repulsive 
force between vortices with anti-parallel orientations.
Due to the similarity in scope and technique, the latter mentioned 
research is the closest to the research presented in chapter \ref{ch:periodic}.



\section{Conclusion}

In this chapter I have discussed my motivation for studying tubes of magnetic 
flux, and the 1-loop \ac{QED} correction to the energy of these configurations. 
In particular, I am interested in the role that flux tubes play in neutron stars 
where the magnetic fields vary rapidly on the Compton wavelength scale. 
In chapter \ref{ch:greensfunc}, I outline progress toward a new Green's 
function technique for computing effective actions for flux tubes, and which 
yields an expression for the quantum-corrected equations of motion for the 
magnetic field. In chapters \ref{ch:WLNumerics} and \ref{ch:WLError}, I outline 
a numerical method implemented on \acp{GPU} to compute the 1-loop effective action 
for magnetic flux tubes. The physics results and conclusions from this method 
are presented in chapter \ref{ch:periodic}.

